\section{Related Work}
\label{sec:related}

\subsection{Robust Multimodal Fusion for Autonomous Perception}
\label{sec:related-fusion}

Camera/LiDAR fusion for 3D detection has moved from simple feature concatenation toward BEV-centric architectures that project both modalities into a shared spatial representation before fusing, with attention-based mechanisms used to weight each modality's contribution per region. A separate line of work has shown that such fusion is not automatically robust: a corrupted or adversarially perturbed single modality can dominate a static or softly-attended fusion decision and collapse detection performance, since nothing in the fusion mechanism verifies that the modalities being combined actually agree geometrically before combining them. The detector this paper monitors adopts a verify-then-fuse design intended to address exactly that failure mode -- a Cross-modal Consistency Probe that computes a per-BEV-cell consistency score between camera and LiDAR features, feeding a Gated Adaptive Fusion Module that quarantines a disagreeing modality at the grid-cell level, trained under an Adversarial Consistency Training objective. We treat this design here strictly as background infrastructure and describe only as much of its mechanism (Section~\ref{sec:method}) as this paper's own monitoring layer depends on, since the monitor itself is detector-agnostic.

\subsection{Conformal Prediction and Anytime-Valid Sequential Testing}
\label{sec:related-conformal}

Split conformal prediction~\cite{VovkGammermanShafer2005,Lei2018JASA} gives distribution-free, finite-sample coverage guarantees for a black-box predictor's outputs by calibrating a nonconformity score on held-out exchangeable data; the guarantee is a batch one, valid at a single fixed evaluation, and does not by itself say anything about monitoring a live, continuously arriving stream. Sequential testing by betting~\cite{Ramdas2023GameTheoretic,WaudbySmithRamdas2024JRSSB} closes that gap: framing a hypothesis test as an accumulating wealth process that is a nonnegative supermartingale under the null lets Ville's inequality bound the false-alarm probability uniformly over all stopping times, giving an anytime-valid guarantee with no fixed monitoring horizon and no post-hoc multiple-testing correction. The direct predecessor for applying this backbone to a live perception-failure problem is Monroy Mu\~noz, Verma, and Timans's work on sequential object-tracking-failure detection~\cite{MonroyMunoz2026WACV}, which frames tracking failure as an e-process over a bounded tracking-quality metric with the betting rate learned adaptively from the metric's own history (approximate GRAPA or Scale-Free Online Gradient Descent), evaluated across four 2D visual-tracking benchmarks. It is a strong, directly applicable baseline, and everywhere in this paper that we compare against a "covariate-blind" betting rule, we mean their rule, reimplemented exactly, not a strawman.

Adjacent recent work conditions conformal guarantees on covariates or system state without adopting the sequential/anytime-valid framing: state-dependent conformal prediction for perception-based autonomous systems~\cite{Geng2025StateDependent} combines conformal error bounds with symbolic reachability analysis but remains a batch procedure; context-aware nonconformity functions for robotic planning and perception~\cite{Kumar2025ContextAware} learn context-conditioned nonconformity scores, with context drawn from learned task representations rather than a cross-modal sensor-disagreement signal specifically; and a value-gated modality-refiner approach~\cite{Liu2026ValueGated} uses cross-modal agreement to gate fusion, but as a learned neural gate rather than a statistical test covariate, and targets sentiment/action-recognition/audio-visual tasks rather than autonomous-vehicle perception. None of these combine sequential anytime-valid testing, spatially-resolved multi-object structure with multiplicity control, and a betting rule conditioned on an external cross-modal disagreement signal -- the combination this paper's monitoring layer targets, following the sharpened contribution statement worked out in the ICRA-targeted predecessor to this paper's monitoring contribution.

\subsection{Conditional Validity and Negative Controls}
\label{sec:related-conditional}

Marginal conformal validity holds on average over the calibration
distribution and can fail badly within identifiable subgroups. Mondrian
conformal prediction~\cite{VovkMondrian2003} addresses this by conditioning
calibration on a taxonomy, giving validity within each category rather than
marginally; the textbook treatment appears in Vovk, Gammerman and
Shafer~\cite{VovkGammermanShafer2005}. The taxonomy is usually a class
label. We apply it with the driving session as the category, which is the
grouping along which we measure exchangeability to fail in real
autonomous-driving data (Section~\ref{sec:exchangeability}), and we report
a case where the resulting benefit reverses.

We are not aware of prior work that reports a null-detector control for a
conformal perception monitor -- that is, rerunning the pipeline with the
monitored model's predictions replaced by a constant to establish that the
monitor responds to model error rather than to dataset structure. The
closest established practice is the use of permutation and label-shuffling
baselines in the broader evaluation literature, which serve the same
epistemic purpose: distinguishing a method's measured effect from
structure already present in the data. Our contribution here is not the
idea of a negative control but the specific, cheap instantiation for
conformal monitoring, together with evidence that it changes conclusions on
real published-style numbers (Section~\ref{sec:null-diagnostic}).

\subsection{Adverse-Weather Datasets for Autonomous Perception}
\label{sec:related-datasets}

Large-scale multimodal driving benchmarks such as KITTI, nuScenes, and Waymo Open Dataset predominantly capture clear-weather conditions, limiting their use for evaluating perception under real distribution shift from adverse weather. Two existing datasets focus specifically on snowy conditions: the Canadian Adverse Driving Conditions dataset (CADC)~\cite{Pitropov2021CADC}, the first fully snowy dataset built for 3D object detection, and the Winter Adverse Driving dataSet (WADS), the first to add 3D semantic segmentation and individual-snowflake labeling; both, however, are reported to have comparatively low-resolution or unimodal sensing relative to more recent captures. Bijelic et al.'s fog- and rain-focused multimodal fusion work~\cite{Bijelic2020SeeingThroughFog} addresses an adjacent adverse-weather regime with a different sensing challenge (fog/rain rather than snow) and is not built around a clear-to-degraded onset transition. Snowy Scenes~\cite{SnowyScenes2026}, described fully in Section~\ref{sec:experiments}, is a substantially more recent (2026) multimodal dataset captured entirely in real snowy conditions with 128-beam LiDAR, RGB, and thermal cameras; the price of that specificity, as we found once real access arrived, is that it contains no clear-weather split at all, which this paper addresses by measuring a real, physically grounded severity ordering across its own weather categories rather than assuming the dataset would hand us the clear/snow contrast originally planned around. This paper's own contribution to that evaluation landscape is not a new dataset, but the first application, to our knowledge, of an anytime-valid sequential monitor to real (rather than synthetically corrupted) multimodal adverse-weather driving data, replicated across two independent such datasets (Snowy Scenes and CADC) rather than one.
